\section{Democratized End-to-End Access to Information}
\label{sec:access}

\subsection{Mission statement}

\Zoo{} treats access to information as a first-class protocol primitive,
not a UX afterthought. The mission, in one sentence:

\begin{quote}\itshape
Every step from user to model to data to output is verifiable,
attributable, and rewards flow to contributors automatically.
\end{quote}

\emph{End-to-end} is the load-bearing word: the property must hold from
the user's keystroke through the model's inference, through the training
data lineage, all the way to the output, and back through the reward
distribution. Skipping any step undermines the chain of custody.

\subsection{Open AI protocol}

\Zoo{} per-\textsc{llm} chains~\cite{zoo-per-llm-chains} commit:

\begin{itemize}[nosep,leftmargin=1.2em]
\item \textbf{Model weights} --- every checkpoint root is Merkle-committed
  on the per-\textsc{llm} chain. Weights are open-weight (downloadable);
  the chain commits the hash, not the bytes.
\item \textbf{Training-data lineage} --- every training shard's hash is
  on the receipts tree. Anyone can audit which examples a model saw.
\item \textbf{Research artifacts} --- every architecture decision,
  benchmark result, and ablation is registered as an on-chain artifact.
\item \textbf{Contribution graph} --- attribution from output back to
  the people whose work fed the model is canonical and on-chain.
\end{itemize}

Compared to closed AI labs (OpenAI, Anthropic), this is a different
ontological commitment. Closed labs treat the weights, the data, and
the architecture as competitive secrets. \Zoo{} treats them as
\emph{public goods} whose provenance is the protocol.

\subsection{Public verifiable inference receipts}

Every inference call to a Hanzo AI Chain endpoint produces a receipt
signed by the \textsc{tee}-attested inference node and committed to the
A-Chain attestation ledger~\cite{lp-134-topology,hanzo-ai-chain}. The
receipt contains:

\begin{itemize}[nosep,leftmargin=1.2em]
\item Hash of the request prompt (or its commitment if private).
\item Hash of the response.
\item Model identifier and chain block height of the model's
  active checkpoint.
\item Wall-clock latency, token count, and compute cost.
\item \textsc{tee} attestation that the response was generated by the
  named model on attested hardware.
\end{itemize}

A user who pays for an inference call has cryptographic proof that the
response was generated by the model they paid for, not a cheaper
substitute. A regulator who audits a model's outputs has a tamper-proof
log of every external use.

\subsection{Zero-knowledge selective disclosure}

The user-side counterpart of public receipts is \emph{selective
disclosure}. A user proves eligibility or capability without revealing
identity:

\begin{itemize}[nosep,leftmargin=1.2em]
\item \emph{``I am over 18''} without revealing date of birth.
\item \emph{``I hold an accredited-investor attestation''} without
  revealing net worth.
\item \emph{``I am a verified researcher with publication credit on
  topic $X$''} without revealing institutional affiliation.
\item \emph{``I have voted in this DAO before''} without revealing
  past votes (consistent with the homomorphic vote tally,
  \S\ref{sec:dao}).
\end{itemize}

These proofs are issued as Groth16 attestations~\cite{groth16}
anchored on Z-Chain and verified on whichever \Lux{} sub-chain
needs the gating decision. The verifier learns only the predicate;
nothing else leaks.

\subsection{Open research registry}

Every \textsc{Zen} model release, every benchmark run, every
architecture decision made by the foundation is committed as an
on-chain research artifact. This includes negative results
(``ablation X did not improve perplexity''), which centralized labs
typically suppress.

The registry is queryable: a researcher can ask, ``which architecture
choices have been tried for problem $Y$?'' and receive a canonical
answer drawn from the on-chain artifacts of every per-\textsc{llm}
chain. This is what we mean by an \emph{open research registry} as
opposed to scattered preprints.

\subsection{Contrast with closed labs}

\begin{longtable}{p{4.5cm}p{4.7cm}p{5.0cm}}
\toprule
\textbf{Property} & \textbf{Closed labs} & \textbf{Zoo open AI protocol} \\
\midrule
\endhead
Model weights & Proprietary, undisclosed & Open-weight; hash committed
   on-chain \\
Training data & Undisclosed; sometimes denied & Lineage hash on
   per-\textsc{llm} receipts tree \\
Inference proof & API trust & TEE-attested receipt on A-Chain \\
Architecture provenance & Marketing post & On-chain artifact registry \\
Reward to contributors & None & Attribution-graph payments
   per per-\textsc{llm} chain \\
Audit by regulator & On request, contractually & On-chain by default \\
\bottomrule
\end{longtable}

\subsection{Democratized end-to-end means}

Putting it together: \emph{democratized end-to-end access to information}
in the \Zoo{} stack means:

\begin{enumerate}[nosep,leftmargin=1.4em]
\item \textbf{User} --- can prove eligibility without revealing
  identity (zero-knowledge selective disclosure).
\item \textbf{Model} --- weights, architecture, and training data are
  on-chain provenanced and open-weight.
\item \textbf{Data} --- every training shard's lineage is auditable;
  every contributor receives attribution.
\item \textbf{Output} --- every inference call produces a verifiable
  receipt; every output can be traced back to the model checkpoint and
  ultimately to the contributors who shaped it.
\item \textbf{Reward} --- royalties and attribution payments flow back
  along the contribution graph automatically, without a central
  treasurer.
\end{enumerate}

This is the operational meaning of \Zoo{}'s mission to be ``an open AI
protocol for research aligned with environment-saving, endangered AI,
and the intersection of AI, blockchain, metaverse, and gaming.''

\subsection{Cross-references}

\begin{itemize}[nosep,leftmargin=1.2em]
\item \cite{zoo-per-llm-chains} --- per-\textsc{llm} chains paper.
\item \cite{hanzo-ai-chain} --- Hanzo AI Chain whitepaper.
\item \Lux{} \textsc{lp-134} --- A-Chain attestation
  topology~\cite{lp-134-topology}.
\item \textsc{zip-0405} --- conservation-aware language models.
\item \textsc{zip-0407} --- decentralized AI training architecture.
\item \textsc{zip-0413} --- foundation language model architecture.
\item \textsc{zip-0606} --- reproducibility attestation.
\end{itemize}
